\chapter{Introduction}
\label{ch:intro}

A supernumerary robotic limb is an extra arm worn on the body. It works
alongside the wearer's own limbs rather than replacing or rehabilitating
them~\cite{yang2021srlreview,prattichizzo2021augmentation}, and the same
idea serves an assistant at a workbench, a surgeon who needs a third hand,
or a person with limited mobility. What makes such a device different from
a fixed industrial robot is simple but consequential: the base of the robot
is a person. It accelerates, sways and steps, and it cannot be told to hold
still.

This thesis reports work on Imperial College London's Dr.\ Octopus platform
(the MUlti-limb Virtual Environment~\cite{muve}), which mounts two Kinova
Gen3 robot arms~\cite{kinova2022gen3}, each with seven joints, on a backpack
frame. The arms are controlled by a \emph{second} person: an operator who
cannot feel the arms, steering them for a wearer who cannot command them.
No previous study has measured, in one system and in the same trial, both
what automated assistance gives such an operator and what it costs the
person carrying the arms (\cref{sec:gap}). Closing that gap is what this
project was built for.

\begin{figure}[htbp]
  \centering
  \includegraphics[width=0.8\linewidth]{photos/platform_scene_hd.png}
  \caption[The platform as built]{The platform as built: two Kinova Gen3
  arms on a backpack frame, worn by a torso mannequin, photographed by the
  room camera. The two teal columns are the laboratory's fixed robots and
  are not part of this system.}
  \label{fig:platform}
\end{figure}

\section{The gap}
\label{sec:gap}

No study measures the operator and the wearer together, under the same
autonomy, in the same trial. Shared autonomy is known to help an operator
when the robot is attached to nobody; apparent autonomy is known to lower a
wearer's trust when the autonomy was a hidden human. Until both are
measured in one system, the exchange rate between what assistance gives the
operator and what it costs the person carrying the arms is unknown. The
full review is in \cref{app:background}.

\section{Aim, objectives and hypotheses}
\label{sec:aim}

The aim was to build and verify, on the platform of \cref{fig:platform}, a
two-person teleoperation system on which that exchange rate can be
measured, and to begin measuring it. Six objectives:

\begin{enumerate}
  \item Design and build an instrumented mannequin controller (the
    \emph{master}) with the robot's joint layout at reduced scale, and
    characterise it against recorded data.
  \item Derive a mapping from master pose to robot command that degrades
    predictably as sensors are lost, and quantify what each loss costs.
  \item Implement a safety architecture suited to a robot mounted beside a
    person, validated by deliberately injecting faults.
  \item Measure the reachable workspace of the mounted arms and which
    manipulation tasks it can support.
  \item Establish a verification method that produces traceable evidence
    for every claim, and apply it throughout.
  \item Run a pilot of direct control against shared-autonomy assistance
    with real operators, and identify what stands between it and the full
    two-person study.
\end{enumerate}

The study the system was built for states five hypotheses, each with a
predicted direction, before any data (\cref{app:study}): task completion
time is shorter under shared autonomy (H1); operator workload is lower
(H2); wearer-reported trust is \emph{lower} under shared autonomy (H3),
opposite to H1 and H2; the wearer's sense of agency falls while ownership
does not (H4); and interference between the two arms falls while the delay
between them does not (H5). The pilot tests H1 directly; no workload
questionnaire was administered, so H2 is untested, and H3--H5 need the
separate wearer the pilot did not have.

The planning report's force-feedback master was not delivered
(\cref{ch:discussionmain,app:development}). \textbf{No result below the
pilot was validated on a physical arm carrying a load.}
